Let \(\mathscr M\) be the locally uniform closure of the moment generating functions of centered laws satisfying \(\lVert\eta\rVert_{\psi_2}\le\psi_{\eta}\) and \(\lvert\kappa_k(\eta)\rvert\ge\delta\).  We first prove that \(\mathscr M\) is compact and that every one of its members retains the hypotheses needed for zero localization.

The sub-Gaussian bound gives a universal \(c_{\psi}<\infty\) such that
\[
 |M(z)|
 \le E[e^{\operatorname{Re}(z)\eta}]
 \le \exp(c_{\psi}\psi_{\eta}^2|z|^2),
 \qquad z\in\mathbb C.
\]
Thus the family is uniformly bounded on each compact disk.  Cauchy's formula on a slightly larger disk makes it equicontinuous on the smaller disk.  Successive finite-net extraction on disks of integer radius, followed by a diagonal argument, therefore gives a locally uniformly convergent subsequence from every sequence of admissible transforms.

To identify every such limit, take corresponding laws.  Their second moments are uniformly bounded, so they are tight.  Pass to a weakly convergent subsequence.  For every real \(t\), the common sub-Gaussian envelope makes \(e^{t\eta}\) uniformly integrable, and hence the transform of the weak limit equals the locally uniform limit on the real axis.  The identity theorem gives equality on \(\mathbb C\).  Centering passes to the limit by uniform integrability.  The sub-Gaussian norm bound is weakly closed: for every \(u>\psi_{\eta}\), Portmanteau's theorem applied to \(x\mapsto e^{x^2/u^2}\), followed by \(u\downarrow\psi_{\eta}\), gives \(\lVert\eta_{\infty}\rVert_{\psi_2}\le\psi_{\eta}\).  Local uniform convergence also gives convergence of all derivatives at zero.  Since the \(k\)-th cumulant is a polynomial in the first \(k\) derivatives of an MGF whose value at zero is one,
\[
 \lvert\kappa_k(M_{\infty})\rvert
 =\lim_m\lvert\kappa_k(M_m)\rvert
 \ge\delta.
\]
Consequently \(\mathscr M\) is compact in the locally uniform topology and every \(M\in\mathscr M\) satisfies the hypotheses of \(\text{lem:zero-localization}\).  In particular, every such \(M\) has a zero in \(\{|z|\le R_0\}\).

Assume first that \(\mathscr M\) is nonempty and fix \(M_{\circ}\in\mathscr M\).  Because \(M_{\circ}(0)=1\), it is not identically zero, and its zeros are isolated.  It consequently has only finitely many zeros in the compact annulus \(R_0\le|z|\le R_1\).  Choose a rational radius
\[
 \rho(M_{\circ})\in(R_0,R_1)
\]
which is not the modulus of any of those zeros.  The circle \(C(M_{\circ})=\{z:|z|=\rho(M_{\circ})\}\) is zero-free, encloses the zero furnished by localization, and satisfies
\[
 m(M_{\circ})
 :=\inf_{z\in C(M_{\circ})}|M_{\circ}(z)|>0,
 \qquad
 N_{C(M_{\circ})}(M_{\circ})\ge1.
\]
The set
\[
 \mathcal U(M_{\circ})
 =\left\{\widetilde M\in\mathscr M:
 \sup_{z\in C(M_{\circ})}
 |\widetilde M(z)-M_{\circ}(z)|<\frac{m(M_{\circ})}{2}
 \right\}
\]
is an open neighborhood in the locally uniform topology.  If \(\widetilde M\in\mathcal U(M_{\circ})\), then
\[
 \inf_{z\in C(M_{\circ})}|\widetilde M(z)|
 \ge\frac{m(M_{\circ})}{2}.
\]
Moreover, \(|\widetilde M-M_{\circ}|<|M_{\circ}|\) on the circle.  Rouch\'e's theorem, equivalently the zero-free straight-line boundary homotopy and the argument principle, gives
\[
 N_{C(M_{\circ})}(\widetilde M)
 =N_{C(M_{\circ})}(M_{\circ})
 \ge1.
\]

The neighborhoods \(\mathcal U(M_{\circ})\) cover the compact set \(\mathscr M\).  Choose a finite subcover centered at \(M_1,\ldots,M_J\) and set
\[
 \rho_j=\rho(M_j),
 \qquad
 a_{\star}=\frac12\min_{1\le j\le J}m(M_j)>0,
 \qquad
 C_j=\{z:|z|=\rho_j\}.
\]
For every \(M\in\mathscr M\), membership in at least one selected neighborhood gives
\[
 N_{C_j}(M)\ge1,
 \qquad
 \inf_{z\in C_j}|M(z)|\ge a_{\star}
\]
for the corresponding index \(j\).  Every treatment-noise transform represented in \(\mathcal P_{\mathrm{NG},n}\) belongs to the original family and hence to \(\mathscr M\), so the conclusion applies uniformly in \(n\).  If the original family is empty, take \(J=1\), any \(\rho_1\in(R_0,R_1)\), and any \(a_{\star}>0\); the covering conclusion is then vacuous.

Only \(k,\delta,\psi_{\eta},R_0,R_1\) enter the existence argument.  The choices of the radii and of the finite subcover are existential.  After one resulting tuple is fixed once and for all it is deterministic and independent of \(n\), \(P\), and the sample, but the argument supplies no effective enumeration, numerical values, or algorithm for finding it.